\section{Ejercicio V}

\subsection{The Big Short (2015)}

En The Big Short, Michael Burry aparece como alguien que entendió antes que los demás que el sistema hipotecario estadounidense estaba completamente mal valuado. Mientras bancos, agencias calificadoras e inversionistas trataban los mortgage-backed securities como activos casi risk-free, Burry hizo algo que parece demasiado básico para Wall Street, leyó los documentos. No se quedó con la etiqueta AAA ni con la narrativa de que las casas siempre suben. Revisó los prospectuses, el collateral y los cash flows, y encontró que muchos bonos supuestamente seguros estaban llenos de hipotecas subprime, borrowers débiles, teaser rates y préstamos que solo funcionaban si el housing market seguía subiendo.

La frase \textquote{I read them} deja en ridículo a todo el sistema. Burry no descubre la burbuja con información privilegiada ni con una fórmula secreta, sino haciendo el trabajo que casi nadie quería hacer. Esa es una de las críticas más fuertes de la película. Wall Street tenía modelos, comités de riesgo, quants, tranches, CDOs, CDS, VaR reports y stress tests, pero el análisis más importante lo hizo alguien leyendo préstamo por préstamo. La sofisticación técnica existía, pero muchas veces funcionaba más como decoración institucional que como verdadera disciplina de riesgo.

Por eso también funciona la escena de \textquote{my quant}. Cuando Danny Moses pregunta si están seguros de las matemáticas, Jared Vennett señala a su especialista y dice \textquote{that's my quant}. Mark Baum responde \textquote{your what?}, y Jared aclara que es su quantitative, su math specialist. La escena es graciosa porque muestra cómo Wall Street vende autoridad técnica sin necesariamente explicar el riesgo. Jared no habla de default correlation, loss given default, prepayment risk ni sensibilidad de los tranches al housing index. Usa al quant como sello de credibilidad. La película no dice que las matemáticas no importen, sino que muchas veces la industria usaba la matemática como marketing.

Burry entiende la diferencia entre rating y riesgo. Un bono podía tener AAA y aun así estar construido sobre una base frágil. El mercado trataba esos mortgage bonds como si la senior tranche fuera casi intocable, pero la estructura dependía de una premisa falsa: que los precios de las casas seguirían subiendo y que los deudores siempre podrían refinanciar. Cuando esa premisa se rompiera, los defaults subirían, los cash flows caerían y las pérdidas empezarían a subir por la estructura del bono.

Por eso el instrumento correcto era el credit default swap. El CDS le permitía comprar protección contra default sobre bonos hipotecarios específicos. Si el bono sobrevivía, Burry pagaba premiums. Si el bono fallaba, el swap pagaba. Era un trade con negative carry, porque tenía que pagar prima cada mes, pero también tenía convexity brutal, downside limitado por las primas y upside enorme si el crédito hipotecario colapsaba. Su frase \textquote{I may have been early, but I'm not wrong} resume exactamente esa tensión. 

La culpa de la crisis del 2008 no está en quienes compraron el short. Burry, Mark Baum, Jared Vennett, Charlie Geller y Jamie Shipley no fabricaron los préstamos NINJA, no obligaron a los brokers a aprobar borrowers sin income, sin job y sin assets, no forzaron a las agencias calificadoras a poner AAA sobre pools llenos de subprime y no crearon la burbuja inmobiliaria. Ellos vieron el mispricing y lo aprovecharon. Eso no destruye el mercado. Eso es price discovery. La pregunta no es por qué alguien shorteó el housing market, sino cómo fue posible que el housing market llegara a ser tan shorteable.

El problema empieza en la originación del crédito. Los mortgage brokers no tenían incentivos para verificar si el borrower podía pagar. Si podían originar un préstamo el viernes y venderlo a un banco el lunes, el riesgo salía de sus manos casi de inmediato. Ese modelo de originate-to-distribute destruyó la disciplina crediticia. El préstamo dejó de ser una evaluación seria de capacidad de pago y se volvió materia prima para la máquina de securitization.

Después vino la ingeniería financiera. Los bancos tomaban miles de hipotecas, las metían en mortgage pools, las convertían en mortgage-backed securities y dividían esos flujos en tranches. La senior tranche cobraba primero y recibía rating alto. Las mezzanine tranches absorbían más riesgo. Las equity tranches quedaban al fondo, con mayor yield y mayor probabilidad de pérdida. En teoría, la estructura podía tener sentido si el riesgo estaba bien medido. El problema es que la diversificación solo funciona si los riesgos no caen todos juntos. En una burbuja nacional de vivienda, esa independencia era una fantasía.

El CDO llevó esa fantasía más lejos. Cuando sobraban tranches BBB difíciles de vender, los bancos las reempaquetaban en collateralized debt obligations. Tomaban pedazos riesgosos de distintos mortgage bonds, los mezclaban y creaban una nueva estructura con nuevas tranches. Por arte de magia, piezas mezzanine podían convertirse otra vez en una supuesta senior tranche con rating alto. Mark Baum lo resume cuando dice que los mortgage bonds son basura y los CDOs son basura envuelta en más basura, recuerda que ningún waterfall, attachment point o modelo de correlación puede salvar una estructura si el collateral está podrido.

El riesgo se volvió todavía más peligroso con los synthetic CDOs. Ahí ya no se empaquetaban solo hipotecas reales, sino exposición sintética mediante credit default swaps sobre los mismos bonos. Una sola hipoteca podía convertirse en muchas capas de exposición financiera. Por eso Vennett explica que, si los mortgage bonds eran el fósforo y los CDOs eran trapos con keroseno, los synthetic CDOs eran la bomba atómica. El daño ya no estaba limitado al tamaño del mercado hipotecario original.

Las agencias calificadoras fueron una pieza central. En teoría, debían estimar default probabilities, expected loss y ratings coherentes. En la práctica, el AAA se volvió parte del packaging. Los bancos necesitaban AAA para vender los productos, los inversionistas necesitaban AAA para justificar la compra y las agencias cobraban de los mismos bancos que pedían la calificación. Así, el supuesto árbitro del riesgo terminó funcionando como proveedor de una etiqueta comercial.

Por eso la quote \textquote{Truth is like poetry. And most people fucking hate poetry} encaja tan bien. La verdad estaba ahí, pero era incómoda. Decía que los préstamos eran malos, que los ratings eran complacientes, que los CDOs reciclaban riesgo y que el mercado inmobiliario no podía subir para siempre. Pero esa verdad dañaba demasiados negocios. Los bancos imprimían fees, los brokers vendían más hipotecas, los ejecutivos cobraban bonuses y los inversionistas obtenían yield extra. Todos tenían incentivos para preferir una mentira rentable antes que una verdad peligrosa.

Mark Baum entiende ese punto desde la rabia. Cuando su equipo escucha a participantes del mercado explicar cómo funcionan los préstamos basura, Baum pregunta por qué están confesando y \textquote{they're not confessing. They're bragging} muestra que el problema ya no era solo técnico, sino moral. La gente no sentía que estaba admitiendo una falla, sentía que estaba presumiendo lo bien que sabía jugar el sistema.

Ben Rickert introduce la parte humana que el lenguaje financiero intenta borrar cuando dice que si ellos tienen razón, la gente pierde casas, empleos, retirement savings y pensions, obliga a mirar el otro lado del trade. Comprar CDS puede ser una posición brillante, pero el evento que lo hace rentable es devastador para millones de personas. No significa que Burry o los demás estén mal por ver la burbuja, sino que no se puede celebrar el payoff como si ocurriera en el vacío.

La quote atribuida a Mark Twain, \textquote{It ain't what you don't know that gets you into trouble. It's what you know for sure that just ain't so}, resume la burbuja. El problema no era ignorancia pura, era falsa certeza. Todos creían saber que las casas siempre subían, que AAA significaba seguridad, que los modelos medían bien el riesgo y que los bancos grandes sabían lo que hacían. Todo eso era precisamente lo que no era cierto.

The Big Short aboga sobre la regulación inteligente. No regulación para matar los mercados, sino para impedir que los mercados se conviertan en casinos apalancados con pérdidas socializadas. Hacían falta reglas más fuertes sobre originación hipotecaria, transparencia en derivados OTC, capital bancario, conflictos de interés en agencias calificadoras, retención de riesgo en securitization y supervisión de productos estructurados.

The Big Short celebra la lucidez de quienes vieron la burbuja antes de que fuera evidente. Burry no destruyó el sistema, era un jugador en el juego. La responsabilidad real está en quienes permitieron que deuda basura se vendiera como AAA, que CDOs reciclaran tranches podridas, que synthetic CDOs multiplicaran exposición fantasma y que la fe en el mercado reemplazara la supervisión.






\subsection{Rogue Trader (1999)}
En \textit{Rogue Trader}, Nick Leeson no quería destruir Barings, quería esconder su pérdida. Fue una persona que comete un error, se avergüenza, intenta cubrirlo y después queda atrapada en una mentira cada vez más grande. Su caída no empieza con una gran decisión criminal, sino con el miedo a admitir que perdió.

Nick llega a Barings con ambición, quiere que lo tomen en serio, quiere demostrar que también puede ocupar un lugar en ese ambiente. Por eso, cuando lo mandan a Asia, ve una oportunidad enorme. En Jakarta demuestra que puede resolver problemas y en Singapur recibe una responsabilidad todavía mayor. Ahí empieza a construirse su imagen de trader exitoso, alguien joven, agresivo y capaz de generar dinero en mercados que otros ni siquiera entienden del todo. Nick empieza a verse como alguien como si tuviera una ventaja especial frente al resto del mercado.

El problema es que esa confianza llega demasiado rápido. Barings le permite manejar tanto las operaciones como el registro administrativo de esas operaciones. Nick no solo toma decisiones de mercado, también tiene la posibilidad de ocultar lo que sale mal.El problema es que ese supuesto control descansa sobre una estructura mal diseñada. Barings le permite manejar el front office y también parte del back office. Nick toma el riesgo y al mismo tiempo puede influir en cómo se registra ese riesgo. Para cualquier institución financiera eso es una falla de control. La persona que ejecuta operaciones no debería ser la misma que confirma, liquida y valida las posiciones. 

Cuando aparece la primera pérdida importante, Nick quiere evitar que Londres se entere, quiere proteger a su equipo y también quiere proteger la imagen que acaba de construir. La cuenta 88888 se vuelve entonces un lugar para guardar lo que no quiere mostrar. Al principio puede parecer una solución temporal, casi como si solo estuviera moviendo un problema para arreglarlo después. Pero esconder una pérdida no la elimina, solo la convierte en una posición más peligrosa, porque ahora el trader ya está gestionando una mentira.

Por eso la frase \textquote{one giant casino} es tan importante. Nick deja de pensar en términos de risk-adjusted return y empieza a pensar como alguien que necesita que la siguiente jugada lo salve. Ya no está preguntando si la posición tiene sentido, si el tamaño es razonable o si el downside está controlado. Está esperando que el mercado se mueva a su favor antes de que alguien descubra el agujero. En ese momento entra en una lógica de doubling down. Pierde, aumenta la posición, vuelve a perder y vuelve a aumentar. Lo que debería haber sido un stop-loss se convierte en una martingala disfrazada de estrategia.

Esa dinámica lo vuelve cada vez más dependiente de la liquidez que recibe desde Londres. Cuando se habla de \textquote{robbing Peter to pay Paul} Nick usa dinero nuevo para cubrir problemas viejos. Pide fondos para supuestas posiciones de clientes, pero en realidad está alimentando una pérdida que pertenece al propio banco. El flujo de efectivo ya no sirve para financiar operaciones sanas, sino para mantener viva una estructura rota. 

La culpa, sin embargo, no puede ponerse solo sobre Nick. Barings también participa en su propia caída porque prefiere creer en el PnL positivo antes que revisar la calidad de ese PnL. Mientras las cifras parecen buenas, sus superiores lo celebran. Les gusta tener a un trader que produce volumen, que domina el piso y que parece estar convirtiendo a Singapur en una máquina de ganancias. Si el dinero entra, el riesgo parece menos urgente. Si el trader gana, su agresividad se interpreta como talento. Si el trader pierde, esa misma agresividad se convierte en irresponsabilidad.

Si Nick hubiera terminado con PnL positivo, probablemente nadie estaría hablando de él como un caso de fracaso. Si sus apuestas hubieran salido bien, tal vez lo habrían presentado como un genio del arbitraje,  alguien capaz de ver oportunidades donde otros solo veían ruido. Quizá habría recibido un bonus enorme, habría vuelto a Londres con estatus de estrella y dentro de Barings lo habrían aplaudido por su agresividad. El sistema muchas veces no castiga el riesgo mal tomado, castiga el riesgo mal tomado cuando explota. Si el trade funciona, se llama visión. Si el trade falla, se llama fraude.

Esa es una idea incómoda porque obliga a pensar en todos los casos que no conocemos. Cuántas historias de éxito pudieron haber nacido de algo parecido, de alguien haciendo doubling down con dinero ajeno, tomando exposición que no entendía del todo y sobreviviendo solo porque el mercado se movió a su favor. Es posible que muchos traders hayan sido premiados por trades que, vistos desde fuera, parecían brillantes, pero que en realidad estaban cargados de riesgo. 

Barings no quiso ver esa posibilidad porque Nick era demasiado útil. Su falta de respeto por las reglas no parecía un problema mientras generaba resultados. Su estilo agresivo era parte de su atractivo. Incluso cuando otros personajes notan que no opera siguiendo los procedimientos normales, el sistema lo defiende porque produce. Esa cultura convierte la imprudencia en virtud. El banco no pregunta con suficiente fuerza quién verifica sus operaciones, por qué necesita tanto funding, qué exposición real tiene abierta o qué pasa si el mercado se mueve violentamente en contra. Se conforma con ver ganancias aparentes, aunque detrás de esas ganancias haya una acumulación de riesgo.

La escena del \textquote{mystery customer X} muestra perfectamente esa distancia entre apariencia y realidad. Nick imagina confesar que el cliente misterioso no existe, que \textquote{we are the customer, Barings}. En los reportes parecía que había clientes externos tomando posiciones, pero en realidad era el propio banco el que estaba expuesto. La contabilidad mostraba una historia cómoda, mientras la economía real decía otra cosa. 

Con el tiempo, Nick deja de tener control sobre sus propias decisiones.  \textquote{The position was running me} resume ese cambio. Al principio él usa el mercado para probar que tiene talento. Después usa el mercado para esconder errores. Finalmente, la posición lo domina. Ya no puede cerrar sin revelar la pérdida, no puede confesar sin destruir su carrera y no puede seguir operando sin aumentar todavía más el riesgo. Está atrapado en una posición con downside creciente y sin salida. En términos de riesgo, su problema ya no es solo direccional. Es un problema de liquidez, de apalancamiento, de reputación y de supervivencia.

El terremoto de Kobe acelera el colapso,, pero no lo causa desde cero. Nick ya estaba demasiado expuesto y Barings ya era demasiado ciego.  \textquote{50 million quid... in one day} no impresiona solo por la cantidad, sino porque muestra el tamaño absurdo del apalancamiento acumulado. Ese es el verdadero blow-up. No solo perder mucho dinero, sino descubrir demasiado tarde que el downside era mucho más grande de lo que todos querían admitir.

Al final, cuando Nick habla de Barings y dice \textquote{wrong sort of person}, parece cerrar el conflicto social de la película. Pero el problema no fue simplemente que Barings contratara a alguien como Nick. El problema fue que quiso aprovechar su ambición, su hambre y su agresividad sin construir controles alrededor de su poder. El banco lo aplaudió mientras parecía imprimir PnL y lo convirtió en culpable absoluto cuando todo salió mal. Nick fue responsable de sus mentiras, de sus falsificaciones y de sus apuestas, pero Barings fue responsable de haber creado un entorno donde esas apuestas podían crecer sin límites claros.

Por eso Rogue Trader no es solo una película sobre finanzas, es una película sobre incentivos. Nick se engaña pensando que puede recuperar todo antes de que alguien lo descubra. Barings se engaña pensando que un PnL positivo prueba que hay control.  La historia de Nick Leeson muestra que una pérdida pequeña puede convertirse en una catástrofe cuando nadie quiere hacer el risk check a tiempo. Y también porque recuerda que el éxito financiero, cuando no se entiende cómo fue producido, puede ser simplemente una posición mal gestionada esperando su momento de explotar.